\begin{definition}[Point-Slope Form]
Let $M = \mathbb{R}^2$ be the Euclidean plane. Given a fixed point $c = (x_1, y_1) \in A$ and a constant slope $m \in \mathbb{R}$, the \textbf{Point-Slope Form} of the line is the set of all points $(x, y)$ satisfying:
\[ y - y_1 = m(x - x_1) \]
\end{definition}

\begin{remark}[Derivation of Point-Slope Form]
The derivation follows from the geometric property of constant slope:
\begin{enumerate}
    \item Let $P_1 = (x_1, y_1)$ be a fixed element of the set $A$.
    \item Let $P = (x, y)$ be an arbitrary point such that $x \neq x_1$.
    \item By the definition of a line in $\mathbb{R}^2$, the slope $m$ between any two distinct points is invariant:
    \[ m = \frac{y - y_1}{x - x_1} \]
    \item To ensure the equation is defined for all $x$ (including $x = x_1$), we apply the algebraic transformation:
    \[ m(x - x_1) = y - y_1 \]
    \item Thus, $y - y_1 = m(x - x_1)$.
\end{enumerate}
\end{remark}


\begin{definition}[Two-Point Form]
Given two distinct points $x, y \in S$ where $x = (x_1, y_1)$ and $y = (x_2, y_2)$ with $x_1 \neq x_2$, the equation of the line passing through both points is:
\[ y - y_1 = \left( \frac{y_2 - y_1}{x_2 - x_1} \right)(x - x_1) \]
\end{definition}


\begin{definition}[General Linear Form]
The \textbf{General Form} (or Implicit Form) of a line in $\mathbb{R}^2$ is defined by the locus of points $(x, y)$ satisfying the linear equation:
\[ Ax + By + C = 0 \]
where $A, B, C \in \mathbb{R}$ and the condition $A^2 + B^2 \neq 0$ ensures the existence of a non-degenerate line.
\end{definition}

\begin{remark}[Geometric Interpretation of Coefficients]
In the context of vector geometry and Vulkan rendering:
\begin{itemize}
    \item The vector $\mathbf{n} = [A, B]^\top$ is the \textbf{normal vector} to the line.
    \item This form is $C^0$ continuous across all possible rotations, including vertical orientations where $B=0$, making it computationally superior to slope-intercept forms for collision detection.
\end{itemize}
\end{remark}


\begin{definition}[Vector Parametric Form]
A line is defined as the image of a mapping $\mathbf{r}: \mathbb{R} \to \mathbb{R}^2$ given by:
\[ \mathbf{r}(t) = \mathbf{r}_0 + t\mathbf{v} \]
where $\mathbf{r}_0$ is a position vector and $\mathbf{v}$ is a direction vector.
\end{definition}

\begin{remark}[Application to Torus Manifolds]
When simulating Brownian motion on a torus $T^2$, the parametric form is typically used with a modulo operator to handle the boundary conditions:
\[ \mathbf{r}(t) = (\mathbf{r}_0 + t\mathbf{v}) \pmod L \]
where $L$ is the periodicity of the torus.
\end{remark}











\begin{definition}[Intercept Form]
Given a line that intersects the $x$-axis at $(a, 0)$ and the $y$-axis at $(0, b)$, where $a, b \neq 0$, the \textbf{Intercept Form} is defined as:
\[ \frac{x}{a} + \frac{y}{b} = 1 \]
\end{definition}

\begin{remark}[Derivation of Intercept Form]
The derivation utilizes the Two-Point Form with specific intercepts:
\begin{enumerate}
    \item Let the two points be $P_1 = (a, 0)$ and $P_2 = (0, b)$.
    \item Calculate the slope $m$:
    \[ m = \frac{b - 0}{0 - a} = -\frac{b}{a} \]
    \item Substitute into the Point-Slope Form using $(a, 0)$:
    \[ y - 0 = -\frac{b}{a}(x - a) \]
    \item Distribute and rearrange the terms:
    \[ y = -\frac{b}{a}x + b \implies \frac{b}{a}x + y = b \]
    \item Divide the entire equation by $b$ to reach the normalized form:
    \[ \frac{x}{a} + \frac{y}{b} = 1 \]
\end{enumerate}
\end{remark}




\begin{remark}[Generalization to $\mathbb{R}^n$ via Affine Combinations]
In an arbitrary $n$-dimensional Euclidean space $\mathbb{R}^n$, a line is not represented by a single equation (unless $n=2$), but rather as an \textbf{affine hull} of two distinct points.

\begin{enumerate}
    \item Let $\mathbf{x_1}, \mathbf{x_2} \in \mathbb{R}^n$ be two distinct vectors representing fixed points in space.
    \item The direction of the line is given by the displacement vector $\mathbf{v} = \mathbf{x_2} - \mathbf{x_1}$.
    \item The "Two-Point Form" in $\mathbb{R}^n$ is expressed parametrically as:
    \[ \mathbf{x}(t) = \mathbf{x_1} + t(\mathbf{x_2} - \mathbf{x_1}) \quad \text{for } t \in \mathbb{R} \]
    \item This can be rewritten in the form of a \textbf{convex/affine combination}:
    \[ \mathbf{x}(t) = (1 - t)\mathbf{x_1} + t\mathbf{x_2} \]
    \item \textbf{Geometric Physicality:} 
    \begin{itemize}
        \item When $t=0$, the "particle" is at $\mathbf{x_1}$.
        \item When $t=1$, the "particle" is at $\mathbf{x_2}$.
        \item For $t \in (0, 1)$, the point lies on the \textbf{line segment} between them. This is the foundation for linear interpolation (LERP) used constantly in shader programming and Vulkan vertex processing.
    \end{itemize}
    \item \textbf{Analytic Requirement:} In $\mathbb{R}^n$, to describe this line implicitly (without a parameter $t$), you would need a system of $n-1$ independent linear equations, which represents the line as the intersection of $n-1$ hyperplanes.
\end{enumerate}
\end{remark}









\begin{definition}[Normal Form / Hesse Normal Form]
A line in $\mathbb{R}^2$ is defined by the length of the perpendicular segment from the origin to the line, $p$, and the angle $\alpha$ that this segment makes with the positive $x$-axis:
\[ x \cos \alpha + y \sin \alpha - p = 0 \]
where $p \ge 0$ is the distance from the origin.
\end{definition}

\begin{remark}[Relation to Vulkan and Distance Functions]
In the context of Signed Distance Fields (SDFs), the expression $f(x,y) = x \cos \alpha + y \sin \alpha - p$ yields the \textbf{signed distance} from any point $(x,y)$ to the line. This is a direct application of the \texttt{Lipschitz} property where the Lipschitz constant $K=1$, making it ideal for GPU-based proximity queries in your pursuit simulation.
\end{remark}

\begin{definition}[Symmetric Form]
Given a point $(x_1, y_1)$ and direction cosines $(\cos \theta, \sin \theta)$, the \textbf{Symmetric Form} is:
\[ \frac{x - x_1}{\cos \theta} = \frac{y - y_1}{\sin \theta} = r \]
where $r$ represents the signed distance from the point $(x_1, y_1)$ to an arbitrary point $(x, y)$ on the line.
\end{definition}


% =========================================================
% Proof: Differentiability Implies Continuity
% Source: [path to notes file relative to project root]
% Mode: standard
% =========================================================

\clearpage
\subsection*{Differentiability Implies Continuity}
\label{prf:differentiability-implies-continuity}

\begin{remark*}[Return]
\hyperref[thm:differentiability-implies-continuity]{$\leftarrow$ Back to Theorem (Differentiability Implies Continuity) in Notes}
\end{remark*}

\begin{remark*}[Claim]
Let $A \subseteq \mathbb{R}$, let $f : A \to \mathbb{R}$, and let $c \in A$ be a limit point of $A$.
Assume that $f$ is differentiable at $c$. Then $f$ is continuous at $c$.
\end{remark*}

\subsubsection*{Professional Standard Proof}

\begin{proof}
Since $f$ is differentiable at $c$, the limit
\[
\lim_{x \to c} \frac{f(x)-f(c)}{x-c}
\]
exists. Let
\[
L:=f'(c).
\]
For $x \in A$ with $x \neq c$,
\[
f(x)-f(c)=\frac{f(x)-f(c)}{x-c}\,(x-c).
\]
As $x \to c$, the first factor tends to $L$, and the second factor tends to $0$. Therefore, by the Product Rule for Limits,
\[
\lim_{x \to c} \bigl(f(x)-f(c)\bigr)=L\cdot 0=0.
\]
Hence
\[
\lim_{x \to c} f(x)=f(c).
\]
Therefore $f$ is continuous at $c$.
\end{proof}

\subsubsection*{Detailed Learning Proof}

\begin{proof}
\textbf{Step 1. Record exactly what differentiability gives.}

Differentiability at $c$ means that the difference quotient has a finite limit:
\[
\lim_{x \to c} \frac{f(x)-f(c)}{x-c}=f'(c).
\]
Set
\[
L:=f'(c).
\]
This is the crucial starting point. Differentiability gives more than continuity: it gives control of the \emph{ratio}
\[
\frac{f(x)-f(c)}{x-c}.
\]
The proof will recover continuity by multiplying this controlled ratio by the small quantity $(x-c)$.

\medskip

\textbf{Step 2. Rewrite the vertical difference in a useful form.}

For every $x \in A$ with $x \neq c$,
\[
f(x)-f(c)=\frac{f(x)-f(c)}{x-c}\,(x-c).
\]
This identity is purely algebraic: the factor $(x-c)$ is inserted and canceled. The point of the rewrite is structural. Continuity concerns the quantity
\[
f(x)-f(c),
\]
while differentiability concerns the quotient
\[
\frac{f(x)-f(c)}{x-c}.
\]
This factorization is the bridge between the two notions.

\medskip

\textbf{Step 3. Identify the limit of each factor.}

From Step 1,
\[
\lim_{x \to c} \frac{f(x)-f(c)}{x-c}=L.
\]
Also,
\[
\lim_{x \to c} (x-c)=0.
\]
The second limit is immediate because subtracting the constant $c$ shifts the variable $x$ but does not change the fact that $x \to c$.

\medskip

\textbf{Step 4. Pass to the limit in the product.}

Using the factorization from Step 2 and the two limits from Step 3, the Product Rule for Limits gives
\[
\lim_{x \to c} \bigl(f(x)-f(c)\bigr)
=
\left(\lim_{x \to c} \frac{f(x)-f(c)}{x-c}\right)
\left(\lim_{x \to c} (x-c)\right)
=
L \cdot 0
=
0.
\]
This is the decisive step. The derivative may converge to any real number $L$, but the factor $(x-c)$ tends to $0$, so the whole product tends to $0$. Thus the vertical change
\[
f(x)-f(c)
\]
vanishes as $x \to c$.

\medskip

\textbf{Step 5. Convert the limit of the difference into continuity.}

The statement
\[
\lim_{x \to c} \bigl(f(x)-f(c)\bigr)=0
\]
is equivalent to
\[
\lim_{x \to c} f(x)=f(c).
\]
But this is exactly the definition of continuity of $f$ at $c$. Therefore $f$ is continuous at $c$.

\medskip

\textbf{Conclusion.}
Differentiability at $c$ forces the difference quotient to have a finite limit. Factoring
\[
f(x)-f(c)
=
\frac{f(x)-f(c)}{x-c}\,(x-c)
\]
then shows that the function value difference tends to $0$ as $x \to c$. Hence
\[
\lim_{x \to c} f(x)=f(c),
\]
so $f$ is continuous at $c$.
\end{proof}

\begin{remark*}[Proof structure]
The argument is a direct factorization proof. Differentiability controls the quotient
\[
\frac{f(x)-f(c)}{x-c},
\]
while continuity requires control of the difference
\[
f(x)-f(c).
\]
The proof inserts the factor $(x-c)$ to convert the derivative datum into the continuity datum. Structurally, the theorem says that differentiability is stronger than continuity because the derivative already contains the continuity estimate in factored form.
\end{remark*}

\begin{remark*}[Dependencies]
\begin{itemize}
  \item \hyperref[def:derivative-at-a-point]{Definition of derivative at a point}: used to obtain the limit of the difference quotient.
  \item \hyperref[def:continuous-at-a-point]{Definition of continuity at a point}: used in the final step to identify the conclusion.
  \item \hyperref[thm:product-rule-for-limits-functions]{Product Rule for Limits}: used to pass from the limits of the two factors to the limit of their product.
  \item \hyperref[thm:algebra-of-limits-functions]{Algebra of Limits for Functions}: used implicitly in passing from
  \(\lim_{x\to c}(f(x)-f(c))=0\) to \(\lim_{x\to c}f(x)=f(c)\), if that equivalence has already been isolated as a theorem.
\end{itemize}
\end{remark*}




\begin{tcolorbox}[colback=gray!6, colframe=gray!40, arc=2pt,
  left=8pt, right=8pt, top=6pt, bottom=6pt,
  title={\small\textbf{Differentiation}},
  fonttitle=\small\bfseries]

\begin{center}
\small
\begin{tabular}{c}
Real Line Foundations $\;\to\;$ Sequences $\;\to\;$ Metric Spaces \\[2pt]
$\to\;$ Limits of Functions $\;\to\;$ Continuity $\;\to\;$ \textbf{Differentiation}
\end{tabular}
\end{center}

\medskip
\noindent\textbf{How we got here.}
The derivative is defined as a limit of a quotient, so the chapter depends
directly on the limit theory of functions and on the pointwise notion of
continuity. The earlier chapters establish the real-number structure,
sequential and metric control, and the algebra of function limits needed to
make difference quotients precise.

\medskip
\noindent\textbf{What this chapter builds.}
This chapter develops differentiation as the first local linear invariant of a
real-valued function. It begins with the definition of the derivative at a
point, establishes that differentiability is stronger than continuity, proves
the basic differentiation rules, and then develops the structural consequences:
local extrema criteria, Rolle's Theorem, the Mean Value Theorem, monotonicity,
and the first global information extracted from local rate of change.

\medskip
\noindent\textbf{Where this leads.}
Differentiation opens the path to the integral, Taylor expansion, and the
analysis of function behavior through approximation by linear and higher-order
models. The Mean Value Theorem, in particular, becomes the bridge from local
derivative information to global control on intervals, and it will later
interact with integration through the Fundamental Theorem of Calculus.

\end{tcolorbox}



% =========================================================
% Definition: Supremum
% =========================================================

\begin{tcolorbox}[colback=defbox, colframe=defborder, arc=2pt,
  left=6pt, right=6pt, top=4pt, bottom=4pt,
  title={\small\textbf{Definition (Supremum)}},
  fonttitle=\small\bfseries]
\begin{definition}
\label{def:supremum}
Let $A \subseteq \mathbb{R}$ and let $s \in \mathbb{R}$. The number $s$ is called the
\emph{supremum} of $A$ if $s$ is an upper bound of $A$ and $s$ is less than or equal to
every upper bound of $A$.
\end{definition}
\end{tcolorbox}

\begin{remark*}[Standard quantified statement]
\[
s = \sup A
\quad\Longleftrightarrow\quad
\bigl(\forall x \in A \; (x \le s)\bigr)
\land
\bigl(\forall u \in \mathbb{R}\; ((\forall x \in A \; (x \le u)) \rightarrow s \le u)\bigr).
\]
\end{remark*}

\begin{remark*}[Predicate statement]
\[
Supremum(s,A)
\quad\Longleftrightarrow\quad
UpperBound(s,A)
\land
\forall u \in \mathbb{R}\; (UpperBound(u,A) \rightarrow s \le u).
\]
\end{remark*}

\begin{remark*}[Negated quantified statement]
\[
s \neq \sup A
\quad\Longleftrightarrow\quad
\neg\Bigl[
\bigl(\forall x \in A \; (x \le s)\bigr)
\land
\bigl(\forall u \in \mathbb{R}\; ((\forall x \in A \; (x \le u)) \rightarrow s \le u)\bigr)
\Bigr].
\]
Equivalently,
\[
\bigl(\exists x \in A \; (x > s)\bigr)
\lor
\bigl(\exists u \in \mathbb{R}\; ((\forall x \in A \; (x \le u)) \land u < s)\bigr).
\]
\end{remark*}





% =========================================================
% Theorem: Every Convergent Sequence Is Bounded
% =========================================================

\begin{theorem}
\label{thm:every-convergent-sequence-is-bounded}
Let $(x_n)$ be a sequence in $\mathbb{R}$ and let $L \in \mathbb{R}$. If $(x_n)$ converges
to $L$, then $(x_n)$ is bounded.
\end{theorem}

\begin{remark*}[Standard quantified statement]
\[
\forall x_n \; \forall L \in \mathbb{R}\;
\Biggl[
\Bigl(
\forall \varepsilon > 0 \; \exists N \in \mathbb{N} \; \forall n \ge N \;
(|x_n - L| < \varepsilon)
\Bigr)
\rightarrow
\Bigl(
\exists M > 0 \; \forall n \in \mathbb{N} \; (|x_n| \le M)
\Bigr)
\Biggr].
\]
\end{remark*}

\begin{remark*}[Predicate statement]
\[
\forall L \in \mathbb{R}\;
\bigl(
ConvergentSequence(x_n,L) \rightarrow BoundedSequence(x_n)
\bigr).
\]
\end{remark*}

\begin{remark*}[Negated quantified statement]
\[
\exists x_n \; \exists L \in \mathbb{R}\;
\Biggl[
\Bigl(
\forall \varepsilon > 0 \; \exists N \in \mathbb{N} \; \forall n \ge N \;
(|x_n - L| < \varepsilon)
\Bigr)
\land
\Bigl(
\forall M > 0 \; \exists n \in \mathbb{N} \; (|x_n| > M)
\Bigr)
\Biggr].
\]

This means there exists a sequence that converges to $L$ but escapes every bound: for every
$M > 0$, one can find a term of the sequence with $|x_n| > M$. In other words, the sequence
is unbounded despite converging.
\end{remark*}

\begin{remark*}[Contrapositive quantified statement]
\[
\forall x_n \; \forall L \in \mathbb{R}\;
\Biggl[
\Bigl(
\forall M > 0 \; \exists n \in \mathbb{N} \; (|x_n| > M)
\Bigr)
\rightarrow
\Bigl(
\exists \varepsilon_0 > 0 \; \forall N \in \mathbb{N} \; \exists n \ge N \;
(|x_n - L| \ge \varepsilon_0)
\Bigr)
\Biggr].
\]

This means that if a sequence is unbounded, then it cannot converge: there is a fixed
$\varepsilon_0 > 0$ such that no matter how far out in the sequence one goes, some term
remains at least $\varepsilon_0$ away from $L$.
\end{remark*}
